% =====================================================================
%  Formulario de inercias, centroides y momentos de inercia de area
%  Vibraciones Mecanicas - Universidad Rafael Landivar
%  Santiago Cruz Rivera - Carne 1310823
%  Compilar con: pdflatex Formulario_Inercias.tex
% =====================================================================
\documentclass[10pt,a4paper]{article}
\usepackage[utf8]{inputenc}
\usepackage[T1]{fontenc}
\usepackage[spanish]{babel}
\usepackage{amsmath,amssymb}
\usepackage[margin=1.2cm]{geometry}
\usepackage{multicol}
\usepackage{array}
\usepackage{booktabs}

\setlength{\parindent}{0pt}
\setlength{\columnsep}{0.9cm}
\pagestyle{empty}

% Compactar el espaciado alrededor de las ecuaciones
\setlength{\abovedisplayskip}{3pt}
\setlength{\belowdisplayskip}{3pt}
\setlength{\abovedisplayshortskip}{2pt}
\setlength{\belowdisplayshortskip}{2pt}

\newcommand{\seccion}[1]{%
  \vspace{2mm}\hrule\vspace{1mm}%
  \textbf{\footnotesize\MakeUppercase{#1}}\vspace{1mm}\hrule\vspace{2mm}}

% Barra de centroide
\newcommand{\bx}{\bar{x}}
\newcommand{\by}{\bar{y}}

\begin{document}
\small

\begin{center}\renewcommand{\arraystretch}{1.45}\normalsize
  \textbf{\large FORMULARIO — INERCIAS, CENTROS DE MASA Y MOMENTOS DE INERCIA}
\end{center}
\vspace{-2mm}

% =====================================================================
\seccion{Momentos de inercia de masa}

Definición general y radio de giro:
\begin{equation*}
  I = \int r^{2}\,dm = \sum_i m_i r_i^{2},
  \qquad
  k = \sqrt{\frac{I}{m}}
  \qquad [\mathrm{kg\cdot m^{2}}]
\end{equation*}

Teorema de ejes paralelos (Steiner). $d$ = distancia entre el eje que pasa
por el centro de masa y el eje paralelo considerado:
\begin{equation*}
  I_{O} = \bar{I}_{cm} + m\,d^{2}
\end{equation*}

\begin{center}\renewcommand{\arraystretch}{1.45}\normalsize
\begin{tabular}{@{}llc@{}}
\toprule
\textbf{Cuerpo} & \textbf{Eje} & \textbf{Momento de inercia} \\
\midrule
Cilindro sólido o disco   & eje de simetría (longitudinal) & $I = \frac{1}{2}MR^{2}$ \\
Cilindro sólido           & diámetro central (transversal) & $I = \frac{1}{4}MR^{2} + \frac{1}{12}ML^{2}$ \\
Anillo (aro delgado)      & eje de simetría                & $I = MR^{2}$ \\
Anillo (aro delgado)      & un diámetro                    & $I = \frac{1}{2}MR^{2}$ \\
Esfera sólida             & un diámetro                    & $I = \frac{2}{5}MR^{2}$ \\
Delgada capa esférica     & un diámetro                    & $I = \frac{2}{3}MR^{2}$ \\
Varilla delgada           & perpendicular, por el centro   & $I = \frac{1}{12}ML^{2}$ \\
Varilla delgada           & perpendicular, por un extremo  & $I = \frac{1}{3}ML^{2}$ \\
\bottomrule
\end{tabular}
\end{center}

\vspace{1mm}
Comprobación de Steiner con la varilla:
$\frac{1}{12}ML^{2} + M\!\left(\frac{L}{2}\right)^{\!2} = \frac{1}{3}ML^{2}$.

% =====================================================================
\seccion{Centro de masa y centroide — fórmula del punto promedio}

El centroide es un \emph{promedio ponderado} de posiciones. El peso de cada
posición es la magnitud que se reparte (masa, área, longitud o volumen):

\begin{multicols}{2}

Sistema discreto de masas (punto promedio):
\begin{equation*}
  \bx = \frac{\displaystyle\sum_i m_i\,x_i}{\displaystyle\sum_i m_i},
  \qquad
  \by = \frac{\displaystyle\sum_i m_i\,y_i}{\displaystyle\sum_i m_i}
\end{equation*}

Cuerpo continuo:
\begin{equation*}
  \bx = \frac{\displaystyle\int x\,dm}{\displaystyle\int dm},
  \qquad
  \by = \frac{\displaystyle\int y\,dm}{\displaystyle\int dm}
\end{equation*}

\columnbreak

Áreas compuestas (el caso más usado):
\begin{equation*}
  \bx = \frac{\displaystyle\sum_i \bx_i\,A_i}{\displaystyle\sum_i A_i},
  \qquad
  \by = \frac{\displaystyle\sum_i \by_i\,A_i}{\displaystyle\sum_i A_i}
\end{equation*}

Área continua:
\begin{equation*}
  \bx = \frac{\displaystyle\int x\,dA}{A},
  \qquad
  \by = \frac{\displaystyle\int y\,dA}{A}
\end{equation*}

\end{multicols}

Los productos $\sum \bx_i A_i$ y $\sum \by_i A_i$ son los \textbf{primeros momentos
de área} $Q_y$ y $Q_x$. Un agujero o recorte entra en la suma con $A_i$
\emph{negativa}. Si el cuerpo es homogéneo, centro de masa y centroide coinciden;
si hay un eje de simetría, el centroide cae sobre él.

\vspace{2mm}
\begin{center}\renewcommand{\arraystretch}{1.45}\normalsize
\begin{tabular}{@{}lccc@{}}
\toprule
\textbf{Figura} & $\bx$ & $\by$ & \textbf{Área} \\
\midrule
Área triangular (base $b$, altura $h$)      & ---                 & $\frac{h}{3}$        & $\frac{bh}{2}$ \\
Área de cuarto de círculo                   & $\frac{4r}{3\pi}$  & $\frac{4r}{3\pi}$    & $\frac{\pi r^{2}}{4}$ \\
Área semicircular                           & $0$                 & $\frac{4r}{3\pi}$    & $\frac{\pi r^{2}}{2}$ \\
Área de cuarto de elipse                    & $\frac{4a}{3\pi}$  & $\frac{4b}{3\pi}$    & $\frac{\pi ab}{4}$ \\
Área semielíptica                           & $0$                 & $\frac{4b}{3\pi}$    & $\frac{\pi ab}{2}$ \\
Área semiparabólica                         & $\frac{3a}{8}$     & $\frac{3h}{5}$       & $\frac{2ah}{3}$ \\
Área parabólica                             & $0$                 & $\frac{3h}{5}$       & $\frac{4ah}{3}$ \\
Enjuta parabólica \quad $y = kx^{2}$        & $\frac{3a}{4}$     & $\frac{3h}{10}$      & $\frac{ah}{3}$ \\
Enjuta general \quad $y = kx^{n}$           & $\frac{n+1}{n+2}\,a$ & $\frac{n+1}{4n+2}\,h$ & $\frac{ah}{n+1}$ \\
Sector circular (semiángulo $\alpha$)       & $\frac{2r\sin\alpha}{3\alpha}$ & $0$      & $\alpha r^{2}$ \\
\bottomrule
\end{tabular}
\end{center}

\vspace{1mm}
En el área triangular, $\by = h/3$ se mide desde la base; $\bx$ queda a un tercio
del camino desde cada lado según la orientación de los ejes. En el sector circular
$\alpha$ va en radianes y se mide desde la bisectriz.

% =====================================================================
\seccion{Momentos de inercia de área}

Definiciones. $I_x$ e $I_y$ son momentos \emph{rectangulares}; $J_O$ es el
momento \emph{polar}:
\begin{equation*}
  I_x = \int y^{2}\,dA,
  \qquad
  I_y = \int x^{2}\,dA,
  \qquad
  J_O = \int r^{2}\,dA = I_x + I_y
  \qquad [\mathrm{m^{4}}]
\end{equation*}

Teorema de ejes paralelos para áreas ($\bar{I}$ = respecto al eje centroidal):
\begin{equation*}
  I = \bar{I} + A\,d^{2},
  \qquad
  J_O = \bar{J}_C + A\,d^{2}
\end{equation*}

Radio de giro de área:
\begin{equation*}
  r_x = \sqrt{\frac{I_x}{A}},
  \qquad
  r_y = \sqrt{\frac{I_y}{A}},
  \qquad
  r_O = \sqrt{\frac{J_O}{A}} = \sqrt{r_x^{2}+r_y^{2}}
\end{equation*}

\begin{center}\renewcommand{\arraystretch}{1.45}\normalsize
\begin{tabular}{@{}lll@{}}
\toprule
\textbf{Figura} & \textbf{Ejes centroidales} & \textbf{Ejes en la base / origen} \\
\midrule
Rectángulo ($b\times h$)
  & $\bar{I}_{x'} = \frac{1}{12}bh^{3}$
  & $I_{x} = \frac{1}{3}bh^{3}$ \\
  & $\bar{I}_{y'} = \frac{1}{12}b^{3}h$
  & $I_{y} = \frac{1}{3}b^{3}h$ \\
  & $\bar{J}_{C} = \frac{1}{12}bh\,(b^{2}+h^{2})$
  & \\
\midrule
Triángulo ($b\times h$)
  & $\bar{I}_{x'} = \frac{1}{36}bh^{3}$
  & $I_{x} = \frac{1}{12}bh^{3}$ \\
\midrule
Círculo (radio $r$)
  & $\bar{I}_{x} = \bar{I}_{y} = \frac{1}{4}\pi r^{4}$
  & $J_{O} = \frac{1}{2}\pi r^{4}$ \\
\midrule
Semicírculo (radio $r$)
  & ---
  & $I_{x} = I_{y} = \frac{1}{8}\pi r^{4}$ \\
  &
  & $J_{O} = \frac{1}{4}\pi r^{4}$ \\
\midrule
Cuarto de círculo (radio $r$)
  & ---
  & $I_{x} = I_{y} = \frac{1}{16}\pi r^{4}$ \\
  &
  & $J_{O} = \frac{1}{8}\pi r^{4}$ \\
\midrule
Elipse (semiejes $a$, $b$)
  & $\bar{I}_{x} = \frac{1}{4}\pi ab^{3}$
  & $J_{O} = \frac{1}{4}\pi ab\,(a^{2}+b^{2})$ \\
  & $\bar{I}_{y} = \frac{1}{4}\pi a^{3}b$
  & \\
\bottomrule
\end{tabular}
\end{center}

\vspace{1mm}
En el semicírculo y el cuarto de círculo los valores tabulados son respecto a los
ejes $x$, $y$ que pasan por el centro $O$ del círculo original, \emph{no} por el
centroide del área. Para llevarlos al centroide hay que restar $A\bar{y}^{2}$ con
$\bar{y} = 4r/3\pi$.

\vspace{2mm}
Secciones huecas: se restan las propiedades del hueco, siempre que ambas figuras
compartan el mismo eje de referencia.
\begin{equation*}
  I = I_{ext} - I_{int},
  \qquad
  A = A_{ext} - A_{int}
\end{equation*}

Formas frecuentes en vibraciones:
\begin{equation*}
  I_{circ} = \frac{\pi d^{4}}{64},
  \qquad
  J_{0,circ} = \frac{\pi d^{4}}{32} = \frac{\pi R^{4}}{2},
  \qquad
  A_{anillo} = \frac{\pi}{4}\left(D^{2}-d^{2}\right)
\end{equation*}

% =====================================================================
\seccion{Nomenclatura}

\begin{multicols}{3}
\begin{tabbing}
\hspace{1.1cm}\=\kill
$M,\,m$ \> masa [kg] \\
$A$ \> área [m$^2$] \\
$R,\,r$ \> radio [m] \\
$D,\,d$ \> diámetro [m] \\
$L,\,l$ \> longitud [m] \\
$b,\,h$ \> base y altura [m] \\
$a,\,b$ \> semiejes de elipse [m] \\
$d$ \> distancia entre ejes [m] \\
$\alpha$ \> semiángulo del sector [rad] \\
$I$ \> inercia de masa [kg$\cdot$m$^2$] \\
$\bar{I}$ \> inercia respecto al centroide \\
$I_x,I_y$ \> inercia de área [m$^4$] \\
$J_O$ \> inercia polar de área [m$^4$] \\
$Q_x,Q_y$ \> primer momento de área [m$^3$] \\
$k$ \> radio de giro de masa [m] \\
$r_x,r_y$ \> radio de giro de área [m] \\
$\bx,\by$ \> coordenadas del centroide [m] \\
\end{tabbing}
\end{multicols}

\end{document}
